\scanpage{69}
\mspage{II.17}
% archive sheet number: 68
% lines 1–4
dans un $E_{\ell'}$, l'une étant semi-simple abélienne, \deleted{\unsure{donc}} \deleted{l'autre est} avec les $\unsure{\rightrightarrows|=}$ pol. car. de Frob, pour presque tous $m\in\operatorname{Spec}R$, alors l'autre représentation est \uline{résoluble},
% lines 4–7, cancelled block including interlinear revisions
\deleted{\unsure{qui} \unsure{a} \unsure{pour} \unsure{obj.} $E_\ell,E_{\ell'}$ \unsure{provenant} $A$ \unsure{par} $T_\ell,T_{\ell'}$, Supposons que \deleted{\unsure{elle}} \unsure{donc} qu'elle soit \inserted{ess. résoluble} \deleted{\unsure{ess.}} aussi abélienne. Alors \unsure{on trouve} que \inserted{et la représentation $\frac12$ simple associée} \unsure{et elle donne} la représentation \uline{sur $\unsure{\mathbb Q|Q}$}}
% lines 8–10
et pour tt $m\in\operatorname{Spec}R$, la repr. \uuline{$\frac12$ simple} \uuline{associée} à $E_{\ell'}$ \unsure{ou|on} \deleted{\unsure{elle|est}} \deleted{\unsure{alors|elle}} est \unsure{le|la} $\unsure{=|\text{même}}$ que celle qui est définie par $E_\ell$. Choisissons alors $\ell$ tel que
